\documentclass[../../main.tex]{subfiles}
\begin{document}

\begin{center}
  {\fontsize{48pt}{58pt}\selectfont \textbf{2005}\normalsize\textbf{年度}}
\end{center}

\begin{table}[htb]
  \centering
  \caption{出題テーマと難易度}
  \begin{tabular}{|c|p{9cm}|c|}
    \hline
    \textbf{問題番号} & \textbf{テーマ} & \textbf{難易度} \\
    \hline
    第1問 & 対数関数の高階導関数と数列の一般項 & \\
    \hline
    第2問 & 複素数の二次変換と絶対値の最大値 & \\
    \hline
    第3問 & 関数の導関数と反復で定めた数列の収束 & \\
    \hline
    第4問 & 整数の二乗と合同式 & \\
    \hline
    第5問 & カードを使ったゲームの勝率 & \\
    \hline
    第6問 & 三つの円柱条件で定まる立体の体積 & \\
    \hline
  \end{tabular}
\end{table}

\begin{table}[htb]
  \centering
  \caption{解答}
  \renewcommand{\arraystretch}{1.4}
  \begin{tabular}{|c|c|p{9cm}|}
    \hline
    \textbf{問題番号} & \textbf{小問} & \textbf{答え} \\
    \hline
    第1問 & (1) & $a_1=1$，$b_1=-1$，
    $a_{n+1}=-(n+1)a_n+b_n$，$b_{n+1}=-(n+1)b_n$ \\
    \cline{2-3}
    & (2) & $a_n=(-1)^{n+1}n!h_n$，$b_n=(-1)^n n!$ \\
    \hline
    第2問 & -- & $w=-\dfrac{25}{16}$ \\
    \hline
    第3問 & (1) & $x>\dfrac{1}{2}$で
    $0\le f'(x)<\dfrac{1}{2}$（証明問題） \\
    \cline{2-3}
    & (2) & $\displaystyle\lim_{n\to\infty}x_n=1$（証明問題） \\
    \hline
    第4問 & -- & $a=625$ \\
    \hline
    第5問 & (1) & 勝率$\left(\dfrac{a}{N}\right)^2$ \\
    \cline{2-3}
    & (2) & 勝率
    $\dfrac{N(N+1)(2N+1)-a(a+1)(2a+1)}{6N^3}$ \\
    \hline
    第6問 & -- & $V=8\left(\sqrt{2}-\dfrac{4}{3}\right)r^3$ \\
    \hline
  \end{tabular}
\end{table}

\paragraph{試験形態}
6問．試験時間は要確認．

\end{document}
